\begin{annexitem}{annex:p7.ect.1}{thm:p07-partial-arithmetic-conservative}{Graphs and program codes}
\begin{proof}
Here is an explicit elimination of partial terms. For each term \(t\)
form an arithmetic graph formula \(G_t(z)\) with a fresh output
variable. For a variable \(x\), use \(z=x\); for a numeral use
its equality with \(z\). For a primitive-recursive function applied
to terms, existentially bind the values of its arguments and assert
its ordinary graph on those values. For \(u\mathbin{\cdot}v\), use
\[
 \exists a,b,s\bigl(G_u(a)\land G_v(b)\land
                      T(a,b,s)\land U(s)=z\bigr).
\]
Fresh variables in this display are renamed at each occurrence.
Translate an atomic equality by existentially binding its two term
values and equating them; translate other atoms in the same strict
way. Commute the translation with the logical connectives and
quantifiers. Structural induction proves functionality of \(G_t\),
using determinacy of finite computations in the application case.

The graph of every primitive-recursive function is provably total
and functional in HA. The proof follows its definition: zero,
successor and projection are immediate; composition substitutes
unique intermediate outputs; primitive recursion codes the finite
list of earlier outputs and uses arithmetic induction to append one
more value. Induction on the list length proves uniqueness. Coding
finite lists can be done arithmetically, as in the computation coding
of \cref{thm:p04-ce-sigma-one}; this assertion does not add sequence
objects to the domain. Replacing the new function symbols by these
graphs therefore eliminates the initial definitional expansion too.

Translate a derivation line by line. Strictness follows from the
existentially asserted argument values. A definedness premise supplies
the value needed for guarded quantifier introduction or elimination;
functionality makes the choice immaterial. Equality substitution uses
equality of these unique values. The application axiom translates
to its own graph definition. An induction instance translates to an
HA induction instance for the translated formula. Thus a proof of
an original-language sentence translates to an HA proof of that
sentence. This also explains why full induction in the expanded
language is harmless here: the translation produces an ordinary
arithmetic formula, without any restriction on its complexity.

For precision, the program coding can be chosen to code the finite
syntax of zero, successor, projections, composition, primitive
recursion, and strict minimization. The finite computation check
records the subcomputations required by each syntax constructor.
Strict minimization at a value records all preceding defined,
nonzero values as well as the zero value found. The universal search
through such finite traces is a partial-recursive evaluator.
Determinacy is proved by induction through two proposed finite
traces, using determinacy of the subcomputations and the least-zero
condition. These are inductions on natural-number codes or lengths,
and are available in HA.

A code with some arguments fixed is constructed by composition with
constant programs for those arguments. The map taking parameters
to that code is primitive recursive, because it only builds finite
syntax; it does not execute the resulting program. Application to a
variable program code is implemented by the fixed universal evaluator.
Accordingly a term built with application, pairing, and arithmetic
operations can be compiled recursively on its syntax. Each constructor
has a finite transformation taking its input computation certificates
to the output certificate, and a terminating output computation
decomposes into those certificates. Induction on the syntax proves
in HA the equal-definedness equations for composition and parameter
substitution. An iterative program is checked by induction on the
number of iterations, adjoining the certificate for the next step.
This proves the program equations used in the subsequent arguments
inside \(\mathsf{HA}^{\ast}\), not merely in the standard model.
The external machine/recursive equivalence of
\cref{thm:p04-tm-partial-recursive} identifies the resulting computable
functions with the handbook's machine computations; the internal
argument here uses the specified finite-code presentation.
\end{proof}
\end{annexitem}

\begin{annexitem}{annex:p7.ect.2}{thm:p07-formal-realizability-lemmas}{Formal realization}
\begin{proof}
All assertions of this paragraph are proved in
\(\mathsf{HA}^{\ast}\). Definedness follows from the leading
conjunct in every clause. Induction on the formula proves substitution
of defined, equal term values into the realization formula. The
normalized tag clause contains only atoms, conjunction and implication;
all other clauses use those operations and universal quantification.
Consequently every realization formula is in \(\mathcal N\).

For \(D\in\mathcal N\), simultaneously prove reflection and
canonical realization. For an atom choose \(c_D=0\); reflection
is its atomic clause. For a conjunction pair the two canonical
terms, and use the projections for reflection. For a universal
formula use \(c_{\forall xD}=\lambda x.c_D\).
The universal assumption makes each application defined with the
required realizing property; conversely, specialize its realizer
at an arbitrary number and apply reflection for the matrix.
For \(D\to E\), choose \(c_{D\to E}=\lambda a.c_E\),
where \(a\) is a new variable. Under the assumption \(D\to E\),
a realizing input for \(D\) gives \(D\) by reflection, hence
\(E\), and hence a defined canonical realizer for \(E\).
The program code is defined without first evaluating \(c_E\).
Conversely, a realizer for \(D\to E\) applied to \(c_D\)
under assumption \(D\) gives a realizer for \(E\); reflection
then gives \(E\). This simultaneous induction is intuitionistic.

For the stated correspondence of antecedent classes, the search for
a witness to an existential atomic primitive-recursive equation is
strict minimization of its decidable test. If a witness exists, a
bounded search up to that witness terminates; the same computation
is a terminating run of the unbounded search. Conversely a terminating
search certificate exhibits a witness. Both implications have HA
proofs. For a partial-term atom, the graph construction in the
preceding proof binds finitely many computation witnesses with a
primitive-recursive check. Pairing these witnesses gives one
existential primitive-recursive equation. Conjunctions, implications
and universal quantifiers preserve provable equivalence, giving the
claimed correspondence. This is why the primitive-recursive
expansion is specified before using the almost-negative terminology.

We next verify formal soundness by induction on a derivation, allowing
realizers as additional variables for its undischarged assumptions.
These checks produce derivations of the realizing assertions, not
merely external semantic facts. An assumption uses its input realizer.
Conjunction introduction pairs two realizers and elimination projects.
Disjunction introduction pairs tag zero or one with the appropriate
payload. For disjunction elimination, inspect the valid numerical
tag and execute the indicated branch proof on the payload; numerical
tag decidability and the two realizing branch implications prove
correctness. Implication introduction compiles the proof program
with the discharged realizer as input; implication elimination applies
that program. An impossible realizer for the false atom justifies
falsity elimination by the base logic.

Universal introduction abstracts the eigenvariable numerically.
Its absence from the undischarged assumptions is precisely what
keeps their realizing conditions unchanged. Universal elimination
executes at the term's value, supplied by its definedness premise.
Existential introduction pairs that value with the matrix realizer.
Existential elimination projects both witness and realizer and executes
the branch proof at those inputs. Its eigenvariable is absent from
the conclusion and from the remaining assumptions. Extra numerical
parameters of a closed theorem can be specialized to zero after the
rule's required substitutions or abstractions. These checks cover
the ordinary intuitionistic natural-deduction rules with the stated
partial-term guards. Equality and strictness axioms have their atomic
content as a realizing condition; equality substitution also applies
to the recursively formed realizing formulas.

Primitive-recursive defining equations and the universal arithmetic
axioms without existential quantification have the canonical realizers
just constructed. For the predecessor-existence axiom, use the total
primitive-recursive predecessor and the zero/positive numerical test.
Their arithmetic equations verify the selected branch and witness.
For the application graph axiom, a realizing input on its atomic side
gives a terminating application. Search for its finite trace with the
stated output and pair that trace with the atomic matrix realizer.
The application axiom itself proves that the search terminates.
In the reverse direction the trace witness proves the atomic equality,
whose realizer is zero. The conjunction in the trace matrix is
realized by pairing zeros. Abstraction over the numerical parameters
gives realizers of its universal closure. This reasoning uses the
axiom as an arithmetic fact in the target theory; it does not assume
the soundness theorem being proved.

It remains to check induction, which need not belong to
\(\mathcal N\). From a realizer of
\(A(0)\land\forall n(A(n)\to A(n+1))\), let \(a\) be
its first component and \(b\) its second. Compile the iteration
\[
 h(0)=a,\qquad
 h(n+1)=(b\mathbin{\cdot}n)\mathbin{\cdot}h(n).
\]
Induction in \(\mathsf{HA}^{\ast}\) on
\(h(n)\downarrow\land h(n)\mathrel{\mathrm r}A(n)\)
proves that it terminates and realizes the matrix at every number.
The compiled index of \(h\) therefore realizes \(\forall nA(n)\).
This verifies the full induction axiom, with the arbitrary formula
and its numerical parameters preserved. The preceding program-code
lemma validates each compilation and each computation equation used
in these rule checks. Formal soundness follows.

Finally let \(u\) realize the premise of a restricted ECT instance.
For each \(x\), under assumption \(D(x)\), put
\[
 q(x)=(u\mathbin{\cdot}x)\mathbin{\cdot}c_D(x),
 \qquad e=\lambda x.\pi_0q(x).
\]
Canonical realization of \(D\) makes \(q(x)\) a defined realizer
of \(\exists yB(x,y)\). Thus \(e\mathbin{\cdot}x\) is its
witness and \(\pi_1q(x)\) realizes \(B(x,e\mathbin{\cdot}x)\).
To realize the ECT conclusion, return the pair of \(e\) and the
index of the following program: on numerical input \(x\), return
an index which, on any realizer \(d\) of \(D(x)\), returns
\(\langle0,\pi_1q(x)\rangle\).
Reflection for \(D\) supplies the assumption needed for termination.
The zero component realizes the true atomic definedness assertion;
the other realizes the matrix. Nested parameter substitution constructs
all these indices without evaluating \(q(x)\) outside \(D(x)\).
Abstraction over \(u\) and any formula parameters supplies the
required base-theory realizer of the full ECT instance.
\end{proof}
\end{annexitem}

\begin{annexitem}{annex:p7.ect.3}{thm:p07-ect-characterization}{Characterization}
\begin{proof}
Work first inside \(\mathsf{HA}^{\ast}+\mathsf{ECT}_0\).
Induct on \(A\). For an atom, truth gives realizer zero and
realization implies truth. For conjunction, pair witnesses in one
direction and project them in the other. For disjunction, its truth
selects one of two tagged witnesses. Conversely its realizing formula
has a valid, decidable tag; the selected branch and the induction
hypothesis give that disjunct. For existential quantification, pair
the numerical witness with its realizing witness; conversely project
both. The definedness conjunct ensures these projections are available.

For a universal formula, its truth and the induction hypothesis give
\(\forall x\exists a(a\mathrel{\mathrm r}A(x))\).
Apply ECT with antecedent \(0=0\) to obtain an index giving a
realizer at each input. It realizes the universal formula. Conversely
specialize a universal realizer at any number and use the induction
hypothesis for the matrix.

For implication, the induction hypotheses turn its truth into
\[
 \forall a\bigl(a\mathrel{\mathrm r}B\to
                       \exists b(b\mathrel{\mathrm r}C)\bigr).
\]
The antecedent \(a\mathrel{\mathrm r}B\) belongs to
\(\mathcal N\), irrespective of the complexity of \(B\).
ECT therefore supplies an index \(e\) taking every such input
\(a\) to a realizer for \(C\). Since the existentially
quantified \(e\) denotes, this is exactly the implication clause.
Conversely, from \(B\) obtain a realizer by the induction
hypothesis, apply the realizing implication, and use the induction
hypothesis for \(C\). This completes the internal equivalence.

For the external forward implication, extend the formal soundness
construction to a proof using ECT: the preceding lemma already supplies
base-theory realizers of all its additional axioms. A proof of a closed
\(A\) thus yields a closed term \(t\) and a base-theory proof
of \(t\mathrel{\mathrm r}A\), hence of \(t\downarrow\).
The standard natural numbers, with actual partial-recursive application,
satisfy every axiom and rule of \(\mathsf{HA}^{\ast}\).
In particular full induction holds in this structure. Consequently
\(t\) really has some numerical value \(n\).

A particular terminating computation has a finite certificate.
Every primitive-recursive check on the numerals in that certificate
is provable by its defining equations and finitely many arithmetic
steps. Using the application graph axiom at each application node
therefore gives a proof of \(t=\overline n\) in the base theory.
Equality substitution yields
\(\overline n\mathrel{\mathrm r}A\).
This argument uses standard-model soundness externally only to ensure
that a finite computation exists. It does not assert an internal
truth predicate or a primitive-recursive bound for finding \(n\).

Conversely, a base proof of
\(\overline n\mathrel{\mathrm r}A\) gives
\(\exists a(a\mathrel{\mathrm r}A)\) in the extension.
The already proved internal equivalence gives \(A\).
Both directions concern closed sentences; substituting a single
standard numeral for an arbitrary open realizer term would not be
valid in general.
\end{proof}
\end{annexitem}

\begin{annexitem}{annex:p7.ect.4}{thm:p07-unrestricted-ect-inconsistent}{A partial classifier diagonal}
\begin{proof}
All reasoning is intuitionistic in the alleged extension. Write
\(K(x)\) for \(x\mathbin{\cdot}x\downarrow\), and put
\[
 D(x)=K(x)\lor\neg K(x),\qquad
 B(x,y)=(y=0\land K(x))\lor(y=1\land\neg K(x)).
\]
By disjunction elimination, \(D(x)\) implies that there exists a
unique \(y\) with \(B(x,y)\). Unrestricted ECT would supply
an index \(f\) defined and correctly classifying each input for
which \(D(x)\) holds. No assertion \(\forall xD(x)\) has
been made or used.

We give the self-reference explicitly. Let \(s(a,b)\) be the
primitive-recursive index of the program \(z\mapsto
 a\mathbin{\cdot}\langle b,z\rangle\). Compile a binary
program \(V_f(a,z)\) which computes
\(f\mathbin{\cdot}s(a,a)\), returns zero if this value is one,
and otherwise continues forever. Its index \(v_f\), in the paired
input convention, is primitive recursive in \(f\). Put
\(q=s(v_f,v_f)\). Its computation equations give
\[
 q\mathbin{\cdot}z\downarrow
 \quad\longleftrightarrow\quad
 f\mathbin{\cdot}q=1,
 \qquad\hbox{hence}\qquad
 K(q)\longleftrightarrow f\mathbin{\cdot}q=1.
\]
The two directions follow by constructing and decomposing its finite
computation certificate. Thus this equivalence is provable in the
base partial arithmetic, not an appeal to a semantic fixed point.

Assume \(K(q)\). Then \(D(q)\), so the classifier condition
forces \(f\mathbin{\cdot}q=0\): the alternative branch would
contradict \(K(q)\). But the displayed equivalence also forces
\(f\mathbin{\cdot}q=1\). Numerical functionality and \(0\ne1\)
give a contradiction. We have therefore proved \(\neg K(q)\).
Now \(D(q)\) again holds, and the classifier condition forces
\(f\mathbin{\cdot}q=1\), hence \(K(q)\), a final
contradiction. Since the premise had a unique output throughout,
restricting unrestricted ECT to unique existence does not help.
\end{proof}
\end{annexitem}
